\onlyinsubfile{\setcounter{section}{0}}
\section{Introduction}\notinsubfile{\label{sec:intro}}

\par A large literature argues that trust matters for aggregate economic outcomes, including long-run growth \cite{Algan2010}. At the micro level, \cite{jbpglg2016} show that economic performance is not simply increasing in trust: income is maximized at intermediate levels. This paper starts from that result and asks whether a similar nonlinearity appears for realized household investment returns.

\par Although the link between trust and economic growth does not need much justificiation, the mechanism is inherent in finance environments since contracting between borrowers and lenders requires a delayed payoff through the \textit{promise of future repayment}. If trust is too low, households may avoid high-return opportunities. This is not a surprising result given the documented equity premium puzzle. Moreover, if trust is too high, investors may under-monitor risk and prices and suffer otherwise avoidable. Frankly, if this mechanism results in a hump-shaped relationship between income and trust, then the pattern should be especially true in a setting which inherenntly requires trust between contracting parties such as returns.

\par  Recent evidence using Norwegian administrative data documents persistent idiosyncratic return differences across households and scale dependence with wealth \cite{aflgdmlp20}. Related work using U.S. panel data also emphasizes a persistent return component and its relevance for inequality and life-cycle outcomes \cite{Daminato2024}. These findings motivate asking what household characteristics can help explain the persistent component.

\par Replicating the Norwegian measurement environment in the U.S. is difficult. U.S. researchers typically combine survey data with partial administrative information, rather than observing a full administrative panel of wealth flows at the population level \cite{Kennickell2017b}. I therefore use the HRS RAND Longitundinal 2022 version of the data, which is unusually useful for this question because it contains the core ingredients needed to construct returns over time: capital income, asset values across waves, debt, and transaction-related information as well as a measure of trust with nonmissing responses from a significant subset of the sample.

\par I verify whether the trust-income hump shape appears in the HRS and then test for a comparable hump shape regarding trust and returns to net wealth. Since trust is observed in one wave, I follow the literature and treat trust as time-invariant in pooled panel specifications \cite{Lesmeister2019}.

\par The main results are as follows. First, I recover the hump-shaped relationship between trust and total income in the HRS. Second, for returns to net wealth, cross-sectional 2022 results are weaker, but the hump-shaped pattern appears in average-return and pooled panel specifications from 2002-2022, including after adding controls for risk exposure. Third, when moving to the fixed-effects decomposition and second-stage regression on time-invariant covariates, trust loses significance, while race and education remain predictive of the persistent component.

\par In extensions, I include financial literacy and alternative trust aggregates tied more directly to financial institutions. Financial literacy (especially the inflation question and literacy score) explains variation in returns in several specifications. For trust in financial institutions, evidence is mixed in cross-section and average specifications, but more favorable in panel settings, with some specifications again suggesting nonlinear trust effects.

\par The rest of the paper proceeds as follows. Section~\ref{sec:litrev} reviews the literature on trust and returns. Section~\ref{sec:data} describes the HRS data and construction of income, wealth, portfolio shares, and return measures. Section~\ref{sec:results} presents the main empirical results for trust, income, and returns, including panel and fixed-effects specifications. Section~\ref{sec:exts} reports extension results with financial literacy and trust in financial institutions.
